\paragraph{The primitive level equals the local conductor.}
Let $f\in S_2(\Gamma_0(N),\chi)$ be the primitive newform obtained
after absorbing the finite character twist, and let $\Pi_f$ be its
unitary cuspidal automorphic representation. The character $\chi$
is arbitrary; its conductor need only divide $N$.
For every prime $p$, write $a(\Pi_{f,p})$ for the local conductor
exponent. Then
\[
 a(\Pi_{f,p})=v_p(N).
\tag{CL1}
\]
Here is the exact published newform input. In
\cite[Theorem~1, p.~302]{casselman}, the local conductor ideal is
the largest ideal $\mathfrak c$ for which the appropriate central
character eigenspace under $\Gamma_0(\mathfrak c)$ is nonzero;
at that ideal the eigenspace is one-dimensional. The remark on
pp.~306--307 identifies this ideal with the exponent in the local
epsilon factor. Theorem~4 on p.~313 gives the global level spaces:
they vanish unless $\mathfrak c(\Pi)\mid M$, and at
$M=\mathfrak c(\Pi)\prod_p p^{b_p}$ their dimension is
$\prod_p(b_p+1)$. The paragraph immediately following that theorem
identifies all primitive newforms with the minimal level
$M=\mathfrak c(\Pi)$. This proves (CL1).

This statement includes nebentypus. Casselman's classical setup on
p.~309 allows every character $\epsilon$ modulo $M$ and uses
$f|_k\gamma=\epsilon(a)^{-1}f$ for
$\gamma=\left(\begin{smallmatrix}a&b\\c&d\end{smallmatrix}\right)
\in\Gamma_0(M)$. Since $ad\equiv1\pmod M$, this is exactly
$\epsilon(d)f$. Thus the usual nebentypus convention is already
covered. No squarefree-level hypothesis occurs.
